\section{Empirical Results \& Discussion}
\label{sec:results}

We structure our empirical analysis around four core research hypotheses evaluated across the primary 32-task frozen benchmark corpus.

\subsection{$H_1$: Semantic Predictability over Frequency Baselines}
We first evaluate whether intent-driven multi-label classification can reliably anticipate task-conditioned constraints on unseen template families compared to an empirical frequency prior. The frequency prior ranks constraint labels by empirical prevalence across the 48 training tasks and selects the top-$k$ most frequent constraints.
On the unseen template family evaluation ($N=32$ from \texttt{family\_c}), LASR at $k=2$ achieves an $F_1$ score of \textbf{0.863} and a \textbf{68.8\% exact bundle match}, compared to \textbf{$F_1 = 0.250$} and 25.0\% exact match for the Frequency Prior baseline. We resampled the 32 evaluation tasks with replacement for $B=2,000$ paired bootstrap replicates, computing percentile bootstrap confidence intervals:
\begin{equation}
\Delta F_1 = +0.612 \quad [95\% \text{ CI: } +0.431, +0.775], \quad p < 0.001
\end{equation}
This strongly supports $H_1$, demonstrating that task text contains predictive signal for constraint pre-allocation across template families.

\begin{table*}[t]
\centering
\caption{Full-Spectrum Systems Evaluation across 32 Development/Validation Tasks ($N=32$, Gemini 3.5 Flash-Lite)}
\label{tab:systems_evaluation}
\small
\begin{tabular}{lcccccccc}
\toprule
\textbf{System Arm} & \textbf{Reactive?} & \textbf{Init $S_0$} & \textbf{Exp. $\Delta S$} & \textbf{1st-Pass (\%)} & \textbf{Struct. (\%)} & \textbf{Sem. (\%)} & \textbf{Mean $TTS$} & \textbf{Total Tokens} \\
\midrule
Full Schema Control & No & 552.0 & 0.0 & 31.2\% & \textbf{93.8\%} & 90.6\% & 2913.2 & \textbf{96,629} \\
$C_I^*$ Baseline ($k=0$) & No & 339.0 & 0.0 & 0.0\% & \textbf{68.8\%} & 65.6\% & 3182.9 & \textbf{102,407} \\
LASR $k=1$ (Static) & No & 361.1 & 0.0 & 15.6\% & \textbf{78.1\%} & 71.9\% & 2410.6 & \textbf{83,754} \\
LASR $k=1$ (Reactive) & Yes & 361.1 & 45.4 & 15.6\% & \textbf{100.0\%} & 93.8\% & 2410.6 & \textbf{77,140} \\
LASR $k=2$ (Reactive) & Yes & 366.8 & 43.8 & 18.8\% & \textbf{100.0\%} & 93.8\% & 2351.5 & \textbf{75,248} \\
Oracle $O(x)$ (Reactive)* & Yes & 359.0 & 48.3 & 15.6\% & \textbf{96.9\%} & 93.8\% & 2325.3 & \textbf{75,649} \\
\bottomrule
\multicolumn{9}{p{0.92\textwidth}}{\footnotesize *\textit{Note}: Oracle $O(x)$ is a model-conditional reference baseline. Mean $TTS$ is averaged over successful trajectories ($N_{\text{succ}}$); Total Tokens accounts for all attempts across all 32 tasks.} \\
\end{tabular}
\end{table*}

\subsection{$H_2$: Invariant Foundation Requirement on the Evaluated Benchmark}
To identify the minimal invariant floor $C_I^*$ meeting the pre-specified 70\% structural-validity threshold ($\tau_{\text{floor}} = 0.70$), we followed a formal sequential design: \textit{pre-specify threshold} $\rightarrow$ \textit{candidate calibration} $\rightarrow$ \textit{freeze $C_I^*$} $\rightarrow$ \textit{final 32-task systems evaluation}.
During candidate calibration on an 8-task calibration probe, we benchmarked three candidate compression tiers prior to task conditioning:
\begin{itemize}
    \item \textbf{Tier 1 (Standard $C_I^*$, 339 tokens)}: 75.0\% structural validity (6/8 tasks).
    \item \textbf{Tier 2 (Compact $C_I$, 308 tokens)}: 62.5\% structural validity (5/8 tasks).
    \item \textbf{Tier 3 (Minimal $C_I$, 284 tokens)}: 50.0\% structural validity (4/8 tasks).
\end{itemize}
Tier 1 ($C_I^* = 339$ tokens) was the smallest tested candidate meeting the 70\% floor and was frozen for all subsequent experiments.
In the subsequent full 32-task systems benchmark (Table~\ref{tab:systems_evaluation}), this frozen foundation alone without task prediction (\texttt{ci\_k0}) yielded 68.8\% structural validity (22/32 tasks), confirming that without task-conditioned augmentation or reactive recovery, structural validity hovers near the 70\% baseline floor. Compressing below 339 tokens triggers progressive degradation due to missing definitions for nested items. We emphasize that 339 tokens represents the smallest tested candidate meeting this threshold, not an unconstrained universal constant.

\subsection{$H_3$: Evidence Consistent with a Causal Contribution of Reactive Recovery}
\begin{table}[htbp]
\centering
\caption{Paired Confidence Intervals and Hypothesis Tests ($N=32$ Development/Validation Tasks)}
\label{tab:paired_bootstrap}
\resizebox{\columnwidth}{!}{%
\begin{tabular}{lcccc}
\toprule
\textbf{Pairwise Contrast} & \textbf{Metric} & \textbf{Point Est.} & \textbf{95\% Bootstrap CI} & \textbf{Significance Test ($p$)} \\
\midrule
LASR Reactive vs Static & $\Delta$ Structural Succ. & \textbf{+21.9 pp} & [9.4, 37.5] pp & Exact one-sided McNemar/binomial ($p = 0.0078125$) \\
& $\Delta$ Semantic Corr. & \textbf{+21.9 pp} & [9.4, 37.5] pp & Exact one-sided McNemar/binomial ($p = 0.0078125$) \\
\midrule
LASR Reactive vs Full & $\Delta$ Structural Succ. & +6.2 pp & [0.0, 15.6] pp & Exact one-sided McNemar/binomial ($p = 0.250$) \\
& $\Delta$ Semantic Corr. & +3.1 pp & [-6.2, 12.5] pp & Exact one-sided McNemar/binomial ($p = 0.500$) \\
\midrule
Workload Token Lift ($k=1$) & Total API Tokens & \textbf{-20.2\%} & -19,489 tokens & Paired Empirical Workload \\
Workload Token Lift ($k=2$) & Total API Tokens & \textbf{-22.1\%} & -21,381 tokens & Paired Empirical Workload \\
\bottomrule
\end{tabular}%
}
\end{table}
The central systems question is whether reactive in-tree recovery provides empirical evidence consistent with a causal contribution over static adaptive exposure. In Table~\ref{tab:systems_evaluation}, we contrast \texttt{ci\_k1\_static} directly with \texttt{ci\_k1\_reactive}:

\noindent\textbf{Structural Reliability}: Final structural success increases from \textbf{78.1\%} to \textbf{100.0\%} ($\Delta\text{Success} = \mathbf{+21.9\text{ pp}}$, 95\% CI: $[+9.4, +37.5]\text{ pp}$, exact one-sided McNemar/binomial test $p = 0.0078125$).

\noindent\textbf{Semantic Correctness}: Task correctness increases from \textbf{71.9\%} to \textbf{93.8\%} ($\Delta\text{Correctness} = \mathbf{+21.9\text{ pp}}$, 95\% CI: $[+9.4, +37.5]\text{ pp}$, exact one-sided McNemar/binomial test $p = 0.0078125$).

\noindent\textbf{Structural Recovery Rate}: All 27 initially failed trajectories in the benchmark were successfully recovered to satisfy structural validation within the three-attempt budget ($27/27 = 100.0\%$).

As broken down in Table~\ref{tab:cluster_breakdown}, the observed lift is concentrated in Cluster 3 (Specialized Bulk Procurement), where static success was only 50.0\% (4/8 tasks) due to missing warehouse enums, while reactive recovery achieved \textbf{100.0\%} (8/8 tasks).

\begin{table}[htbp]
\centering
\caption{Task Archetype Breakdown: Static Vulnerability vs. Reactive Recovery ($N=8$ per Cluster)}
\label{tab:cluster_breakdown}
\resizebox{\columnwidth}{!}{%
\begin{tabular}{lcccc}
\toprule
\textbf{Task Cluster Archetype} & \textbf{Full (\%)} & \textbf{Static $k=1$} & \textbf{Reactive $k=1$} & \textbf{Observed Lift} \\
\midrule
C1: Routine Regional & 87.5\% & 100.0\% & \textbf{100.0\%} & \textbf{0.0 pp} \\
C2: Expedited Emergency & 100.0\% & 87.5\% & \textbf{100.0\%} & \textbf{+12.5 pp} \\
C3: Specialized Bulk & 87.5\% & 50.0\% & \textbf{100.0\%} & \textbf{+50.0 pp} \\
C4: Multi-Item Depot & 100.0\% & 75.0\% & \textbf{100.0\%} & \textbf{+25.0 pp} \\
\bottomrule
\end{tabular}%
}
\end{table}

\subsection{$H_4$: Total Workload Token Efficiency}
Table~\ref{tab:systems_evaluation} highlights the macro token economics of adaptive resolution across the 32-task workload: Full Schema Control consumed \textbf{96,629 total API tokens} with a mean $TTS$ of 2,913.2 tokens, while reactive LASR cut consumption to \textbf{77,140 tokens} at $k=1$ (\textbf{-20.2\%}) and \textbf{75,248 tokens} at $k=2$ (\textbf{-22.1\% net token reduction}).

\textbf{Observed Correctness vs. Full Schema}: In this benchmark, reactive LASR achieved 93.8\% semantic correctness compared to 90.6\% for Full Schema (+3.1 pp, corresponding to 30/32 vs. 29/32 tasks; unrounded difference +3.125 pp). Because LLM outputs exhibit inherent stochasticity, this difference reflects an observed empirical result on this benchmark distribution rather than evidence that schema reduction intrinsically enhances underlying semantic comprehension. The core claim rests on reliability preservation and workload token reduction.

\textbf{Why Reactive LASR Outperforms Static LASR in Total Tokens}: In \texttt{ci\_k1\_static}, failing tasks cycle through 3 wasted attempts without converging, accumulating 83,754 tokens. In \texttt{ci\_k1\_reactive}, because in-tree schema expansion resolves the root cause immediately on Attempt 2, the system terminates the retry cascade early, saving 6,614 tokens (-7.9\%) relative to static LASR.

\subsection{Emergence of $k=2$ as the Best Observed Operating Point}
While $k=1$ provides the leanest initial schema ($S_0 = 361.1$ tokens), \texttt{ci\_k2\_reactive} ($S_0 = 366.8$ tokens) achieved the lowest total token consumption overall (75,248 tokens) while achieving 100\% structural validity and 93.8\% semantic correctness. By predicting slightly wider constraint bundles ($k=2$) upfront, the system avoids initial failures on Cluster 2 without adding meaningful schema bloat.

\subsection{Decomposing Token Savings: Schema Reduction vs. Retry Avoidance}
A natural systems question is: \textit{How much of the total token savings stems from a leaner schema footprint versus avoiding failed retry turns?} Auditing granular token accounts across the 32 tasks reveals the exact operational balance:

\noindent\textbf{Initial Schema Savings vs. Expansion Cost}: In \texttt{ci\_k2\_reactive}, total schema serialization consumed 24,972 tokens compared to 34,776 for Full Schema—a net saving of \textbf{9,804 schema tokens (-28.2\%)}, even after paying 1,402 tokens of amortized reactive expansion across all attempts ($43.8$ tokens/task).

\noindent\textbf{Retry-Avoidance Savings}: In \texttt{ci\_k1\_static}, failing tasks cycle through 3 attempts without converging, accumulating 75 total attempts across the workload. In \texttt{ci\_k1\_reactive}, in-tree schema expansion resolves omitted rules immediately on Attempt 2, reducing workload attempts to 66 (avoiding 9 full failure rounds). This retry avoidance saves \textbf{4,059 prompt tokens} and \textbf{1,294 completion tokens} (5,353 tokens), plus \textbf{1,261 avoided schema tokens}, yielding 6,614 net API tokens saved (-7.9\%) over the static arm.

\noindent\textbf{Net Workload Impact}: Comparing \texttt{ci\_k2\_reactive} against Full Schema, $45.9\%$ of the 21,381 saved tokens ($9,804$ tokens) derives from leaner schema serialization, while $54.1\%$ ($11,577$ tokens) derives from prompt and completion token savings during multi-turn interactions.

\subsection{Auditable Trajectory Case Study}
Figure~\ref{fig:case_study} details the step-by-step recovery trajectory for task \texttt{BLIND-C3-013}. On Attempt 1, exposed to $C_I^* + \text{top-1 LASR rule}$ (347 tokens), the model generates unconstrained enum \texttt{'WH-EAST-01'}. The schema validator flags three concurrent violations (\texttt{warehouse\_code.enum}, \texttt{order\_priority.enum}, \texttt{supplier\_id.pattern}). LASR intercepts these error paths and recomposes the schema tree in-tree ($347 \rightarrow 419$ tokens). On subsequent attempts, the agent resolves all constraints, converging to 100\% structural and semantic correctness.

\begin{figure}[t]
\centering
\fbox{
\parbox{0.92\columnwidth}{
\scriptsize
\textbf{Task}: \texttt{BLIND-C3-013} (Specialized Bulk) \quad \textbf{Arm}: \texttt{ci\_k1\_reactive} \\
\rule{0.92\columnwidth}{0.4pt}
\textbf{[Attempt 1]} \\
$\bullet$ Exposed Schema: $C_I^* \cup \{\text{items.maxItems}\}$ (347 tokens) \\
$\bullet$ Model Output: \texttt{warehouse\_code = 'WH-EAST-01'} \\
$\bullet$ Validator: FAIL (3 concurrent violations in pass):\\
\quad 1. \texttt{'WH-EAST-01' not in ['WH-US-EAST', \dots]} \\
\quad 2. \texttt{'standard' not in ['STANDARD', \dots]} \\
\quad 3. \texttt{'SUP-ACME1234-US' does not match regex} \\
$\bullet$ Extracted: $\{\texttt{warehouse.enum}, \texttt{priority.enum}, \texttt{supplier.pattern}\}$ \\
\textbf{[In-Tree Recomposition]} Spliced Schema: $347 \rightarrow 419$ tok (+72 tok) \\
\textbf{[Attempt 2]} Model Output: \texttt{'WH-US-EAST'} (Enum Resolved) \\
\textbf{[Attempt 3]} Output: Valid purchase order JSON \\
$\bullet$ Validator: \textbf{PASS} (Structural: 100\%, Semantic: 100\%) \\
$\bullet$ Outcome: \textbf{RECOVERED} (3 attempts, $TTS = 3,626$ tok)
}
}
\caption{Trajectory-level evidence of pass-level multi-error extraction and native in-tree schema recomposition on family\_c evaluation task \texttt{BLIND-C3-013}.}
\label{fig:case_study}
\end{figure}

